% !TEX root = ../main.tex
\section{주장과 증거를 잇는 지도}
\label{ch:appendix-claims}

이 부록은 각 가설(H), 연구 질문(RQ), 그리고 실제 자료로 확인한 주장(C)을 \ref{ch:results}장에 있는 관련 증거와 이어 줘요. 그리고 그 증거가 보여 주지 않는 것도 적어요. 여기서 말하는 증거는 수정 분석(등록 뒤에 금액을 달러로 바꿔 다시 한 분석)의 증거예요. 수정 분석에서는 다섯 토큰의 금액을 달러(USD)로 쟀어요. 이 방법은 등록 분석(미리 등록한 그대로 한 분석)의 결과를 안 뒤에 정했어요(\ref{sec:amount-revision}절). 등록 분석의 결과는 부록~\ref{sec:map-registered}에 남겨 두었어요.

\dissertationSubheading[sec:map-h1]{H1 / RQ1: 순위 갈라짐}

RQ1은 엔도스랭크(EndorseRank)와 AWP(적응형 가중 페이지랭크)가 같은 스마트 계약(블록체인 위에서 저절로 실행되는 프로그램)들을 서로 다른 순서로 세우는지 물어요.

\begin{itemize}
\item 진술: 매칭 코호트(모든 점수가 똑같이 점수를 매기는 계약 묶음)에서 엔도스랭크와 AWP는 계약들에 서로 다른 순위를 매겨요.
\item 증거: \ref{sec:rank-divergence}절, 표~\ref{tab:tie-sensitivity}, 그리고 그림~\ref{fig:rank-scatter}과~\ref{fig:score-distributions}에 있어요. 그 표의 마지막 행은 두 방법 사이의 켄달의 타우 $\tau_b$(두 줄 세우기가 얼마나 같은 순서인지 나타내는 수)와 그 구간(믿을 만한 범위)을 보여 줘요. 각 그래프(점과 화살표로 이은 그림) 밖에 있는 계약을 외톨이 점(화살표가 하나도 없는 점)으로 남겨 둔 경우의 값도 함께 보여 줘요.
\item 해석: 두 순위는 같은 것이 아니에요. 두 점수가 모두 순서를 매기는 쌍 가운데 많은 몫이 서로 반대 순서로 놓여요. 그리고 그래프 밖의 계약을 다루는 규칙은 두 점수의 일치를 바꾸지 않아요. 목표 O1의 기준은 채워졌어요(표~\ref{tab:objectives-achievement}). 등록 분석에서도 그랬어요(표~\ref{tab:registered-vs-revised}).
\end{itemize}

\dissertationSubheading[sec:map-h2]{H2 / RQ2: 같은 기간 일치도}

RQ2는 각 점수가 같은 기간의 허락(allowance, 토큰을 얼마까지 대신 써도 된다는 허락) 점검과 얼마나 가깝게 맞는지 물어요. 그리고 같은 기간의 송금(transfer) 점검과 얼마나 가깝게 맞는지도 물어요.

\begin{itemize}
\item 진술: 같은 기간 안에서, 각 점수는 자기가 걷는 화살표(엣지) 종류의 대리 지표(대신 재는 값)와 가장 잘 맞아요. 엔도스랭크는 허락 지표 가족(비슷한 지표 묶음)과, AWP는 송금 지표 가족과 가장 잘 맞아요. 그리고 C-PR(결합 페이지랭크, 두 그래프를 함께 걷는 점수)의 혼합 점수들은 모든 가족에서 자기의 두 층 사이에 있어요. 금액을 달러로 재면, AWP는 계약에 들어오는 송금 금액의 합을 따라가요. 그 가까운 정도는, 그 송금을 보낸 쪽(송신자)의 수, 곧 들어오는 화살표 수(입차수)를 따라가는 정도와 거의 같아요.
\item 증거: 모든 점수의 지표 가족 평균과 그 구간은 표~\ref{tab:alignment-hybrid}에 있어요. 송금 대리 지표 하나하나의 값은 표~\ref{tab:alignment-transfer}에 있어요. C-PR과 S-PR(보증 씨앗 페이지랭크)의 짝지은 대비는 표~\ref{tab:tau-diff-hybrid}에 있어요. 그래프 밖의 계약을 외톨이 점으로 남겨 둔 같은 기간 계수는 표~\ref{tab:tie-sensitivity}에 있어요.
\item 참고: 이것은 의도한 구성개념(점수가 재려고 하는 생각) 점검이에요. 대리 지표가 그래프와 같은 이벤트(계약이 남기는 기록)에서 나오기 때문에, 일부는 만든 방식 때문에 저절로 그렇게 돼요. 엔도스랭크는 코호트 안에 동점을 조금만 남겨요. 그리고 그래프 밖의 계약을 다루는 규칙은 어떤 계수도 바꾸지 않아요. 계획에서 정한 것은 규칙~A의 같은 기간 부분뿐이에요. 이 부분은 C-PR에 관한 것이에요(\ref{sec:coupled-results}절). 그런데 여기서는 이 부분을 등록 분석의 결과를 안 뒤에 정의한 점수에 적용해요. 그래서 이 비교들 가운데 어느 것도 자기 자료보다 먼저 정해지지 않았어요. 엔도스랭크와 AWP 사이의 짝지은 차이와 대리 지표별 결과는 부록~\ref{ch:appendix-er-awp}에 있어요.
\end{itemize}

\dissertationSubheading[sec:map-h3]{H3 / RQ3: 계산 비용}

RQ3은 같은 계약 집합에서 각 점수를 계산하는 데 비용이 얼마나 드는지 물어요.

\begin{itemize}
\item 진술: 같은 계약들에서, 풀이 프로그램(solver)을 한 번 부를 때 걸리는 시간은 그 그래프의 화살표 수에 반복(한 바퀴) 횟수를 곱한 값을 따라가요. 그래서 더 듬성듬성한(화살표가 적은) 그래프를 쓰는 엔도스랭크가 시간이 가장 적게 걸려요. 반면 C-PR은 두 층을 합친 더 큰 그래프에서 돌아요. 그리고 S-PR은 다시 출발(재시작)을 위해 허락 층을 한 번 더 풀어요. 그래서 C-PR과 S-PR은 AWP보다 오래 걸려요.
\item 증거: 표~\ref{tab:benchmark-runtime}(실행 시간, 표준편차(SD), 메모리, 반복 횟수, $|V|$, $|E|$), 표~\ref{tab:benchmark-scaling}, 그리고 표~\ref{tab:benchmark-runs}의 모든 시간 잰 실행에 있어요.
\item 참고: 모든 점수는 같은 풀이 프로그램을 써요. 그래서 속도는 $|E|$의 성질이지, 어떤 알고리즘이 가진 장점이 아니에요(부록~\ref{sec:pr-sparsity}). 규모에 따른 변화 단계는 다섯 토큰의 코호트 그래프에서 뽑은 부분 표본이에요. 그래서 이 주장은 그 그래프보다 크지 않은 그래프에만 들어맞아요.
\end{itemize}

\dissertationSubheading[sec:map-rq4]{RQ4: 시간 홀드아웃}

RQ4는 $t_1$에 동결한(고정한) 점수가 스펜더 계약(허락을 받아 대신 쓰는 계약)에 나중에 새로 들어오는 것들을 잘 줄 세우는지 물어요. 그리고 어떤 페이지랭크(PageRank)가 자기가 매끄럽게 다듬는(평활하는) 원시 차수(화살표를 그냥 센 수)보다 이것들을 더 잘 줄 세우는지도 물어요.

\begin{itemize}
\item 진술: 각 라벨(나중에 채점에 쓰는 정답 값) 기간이 시작할 때 점수를 동결했어요(W0은 $t_1$, W1은 $T_{\mathrm{obs}}$). 그러면 각 점수는 스펜더 코호트에서 자기 화살표 종류로 나중에 새로 들어온 것들을 가장 잘 줄 세워요. 동결일의 원시 차수와 견주면 두 걷기는 서로 달라요. AWP는 두 기간 모두에서 나중의 새 송금자(새로 송금을 보내온 주소)를 송금 입차수(송금을 보내온 서로 다른 주소 수)가 줄 세우는 것보다 더 잘 줄 세워요. 반면 계획에서 정한 고르게 다시 출발하기(균등 재시작)를 그대로 쓰는 엔도스랭크는, 두 기간 모두에서 나중의 새 허락 쌍(새로 생긴 주인–스펜더 허락)을 받은 허락 수(허락 입차수)가 줄 세우는 것보다 못하게 줄 세워요. 등록 분석도 이 대비들에 같은 부호를 주었어요(표~\ref{tab:registered-vs-revised}).
\item 증거: 표~\ref{tab:holdout-spenders}과~\ref{tab:holdout-diff}(\ref{sec:holdout-results}절)에 있어요. 기간 W1은 표~\ref{tab:fresh-posthoc}에 있어요. 동결 전에 송금을 받은 적이 있는 스펜더는 표~\ref{tab:holdout-receiving}에 있어요. 활동 재시작(활발한 주소에서 더 자주 다시 출발하기)을 쓰는 엔도스랭크는 표~\ref{tab:endorserank-restarts}에 있어요.
\item 주의할 점: 두 라벨은 나중 기간에 관찰한 구성개념이에요. 곧 다섯 토큰에서 새 허락 쌍의 수와 새 송금자의 수예요. 신용이나 거래의 결과가 아니에요. 등록 분석에서는 규칙~A의 두 대비를 계획에서 정했어요. 그리고 규칙~B의 대비는 F4--F6과 함께 W1 등록 문서에서 정했어요. 둘 다 이 연구가 그 판정에 쓸 자료를 뽑기 전에 정했어요. 하지만 고친 점수는 등록 분석의 결과를 안 뒤에 정의했어요(\ref{sec:amount-revision}절). 그래서 고친 점수에서는 어느 대비도 미리 정한 것이 아니에요. 다음 값들은 라벨을 본 뒤에 계산했어요. 엔도스랭크와 AWP의 차수 대비, C-PR ($\lambda=1$)이 받은 허락 수보다 더 얻은 몫, S-PR의 대비, 그리고 엔도스랭크와 AWP 사이의 차이(부록~\ref{ch:appendix-er-awp})예요. AWP의 활동 재시작을 쓰는 엔도스랭크는 두 기간 모두에서 새 허락 쌍을 받은 허락 수가 줄 세우는 것보다 더 잘 줄 세워요. 하지만 그 비교도 라벨을 본 뒤에 했어요. 그래서 그것은 발견이 아니라, 등록해서 시험할 가설이에요.
\end{itemize}

\dissertationSubheading[sec:map-rq5]{RQ5: 결합 연산자}

RQ5는 C-PR이 송금만 걷는 걷기보다 나은지 물어요. 그리고 C-PR이 어느 한 층만 걷는 걷기보다 계약들을 더 잘 줄 세우는지도 물어요.

\begin{itemize}
\item 진술: 고친 점수에서, 두 화살표 종류를 한 번에 걷는 걷기는 한 층만 걷는 것을 이기지 못해요. 규칙~A는 실패해요. C-PR은 새 허락 쌍에서는 허락 걷기보다 나쁘지 않아요. 하지만 새 송금자에서는 송금 걷기보다 나빠요. 그리고 규칙~A의 같은 기간 부분도 실패해요. 규칙~B의 기준도 채워지지 않아요. C-PR은 새 허락 쌍에서 송금 걷기보다 나아요(F1). 그래서 W0에서 본 이득이 W1에서도 다시 나타나요. 하지만 새 송금자에서는 허용 폭을 채우지 못해요(F2). 그리고 트레이더(거래하는 사람)의 청산 없이 닫은 비율(청산, 곧 손실이 너무 커져 거래가 강제로 끝나는 일 없이 거래를 닫은 비율)에서도 아슬아슬하게 채우지 못해요(F3). 여기서는 두 점 추정값이 같지만, 구간의 아래 끝이 $-\delta$보다 아래에 있어요. 첫 번째 민감도 분석(조건을 바꿔 결과가 얼마나 움직이는지 보기)은 AWP를 비교 대상으로 써요. 여기서는 F2가 분명하게 실패해요. 라벨을 본 뒤에 계산한 W0 대비에서, S-PR은 새 송금자에서는 송금 걷기와 같은 수준에 머물러요. 그리고 새 허락 쌍에서는 허락 걷기보다 아래예요. 곧, 허락 화살표는 걷기가 그 화살표를 따라 움직일 때에만 도움이 돼요. 등록 분석도 같은 판정에 이르렀어요(부록~\ref{sec:map-registered}).
\item 증거, 규칙~A: 주 부분은 표~\ref{tab:holdout-diff}에 있어요. 보조 부분은 표~\ref{tab:alignment-hybrid}과~\ref{tab:tau-diff-hybrid}에 있어요(\ref{sec:coupled-results}절).
\item 증거, 규칙~B: 표~\ref{tab:fresh-contrasts}에 있어요. 이 표에는 F1--F3이 있고, 그 옆에 F4--F6과 두 민감도 분석도 함께 보고해요(\ref{sub:fresh-results}절). W0의 탐색적(미리 정하지 않고 살펴본) 대비는 표~\ref{tab:holdout-diff}에 있어요. 등록 분석의 판정과 수정 분석의 판정을 나란히 놓은 것은 표~\ref{tab:registered-vs-revised}에 있어요.
\item 증거, 비용: 표~\ref{tab:benchmark-runtime}의 C-PR 행과 S-PR 행에 있어요.
\item 주의할 점: 이 시험은 다음 조건에서 했어요. $\lambda\in\{0.25,0.5,0.75\}$로 고정했어요. 층은 다섯 토큰의 금액을 달러로 쟀고, 허락 하나하나는 상한 $U$(이보다 큰 금액은 이 값으로 자름)에서 잘랐어요. 순간 이동(다시 출발)은 고르게 했어요. 섞기는 이동 행렬(전이 행렬) 수준에서 했어요. 그리고 체인은 하나만 썼어요. 규칙~A와 $\lambda$ 값들은 계획과 설정 파일에서 정했어요. 계획과 설정 파일은 이 연구의 로그를 하나도 뽑기 전인 2026년 10월 9일에 \texttt{43d8045}로 커밋했어요(\ref{sec:holdout-protocol}절). 규칙~B, 그 허용 폭, 그리고 그 두 민감도 분석은 W0의 결과를 안 뒤에 썼어요. 그리고 스캔~B 전에 \texttt{0d3a057}로 등록했어요(\ref{sub:fresh-holdout}절). 달러로 잰 층과 상한은 등록 분석의 결과를 안 뒤인 2026년 10월 10일에 정했어요(\ref{sec:amount-revision}절). 그래서 위의 판정들은 이 규칙들을 고친 점수에 재분석으로 적용한 것이에요. 등록한 시험이 아니에요.
\end{itemize}

\dissertationSubheading[sec:map-c1c4]{주장 C1--C5}

\ref{ch:results}장은 주장 C1--C6을 말해요. 처음 다섯 주장은 아래의 증거에 기대요. C6의 증거는 부록~\ref{sec:map-c6}에 이어서 나와요.

\begin{itemize}
\item C1: H3/RQ3 아래에 적은 실행 시간 증거예요. 화살표 수의 비와 시간의 비는 \ref{sec:benchmark-results}절과 부록~\ref{sec:pr-sparsity}에 있어요. 속도는 알고리즘이 아니라 더 듬성듬성한 그래프 덕분이에요. 그리고 C-PR과 S-PR은 AWP보다 비용이 더 들어요.
\item C2: 모든 점수가 자기 화살표로 만든 대리 지표와 같은 기간 안에서 보이는 일치예요(표~\ref{tab:alignment-hybrid}). 이것은 의도한 구성개념 점검이고, 일부는 만든 방식 때문에 저절로 그렇게 돼요. AWP는 들어오는 송금의 달러 금액을, 그 송금을 보낸 쪽의 수를 따라가는 것과 거의 같은 정도로 가깝게 따라가요(표~\ref{tab:alignment-transfer}). 동점과, 그래프 밖의 계약을 다루는 규칙은 표~\ref{tab:tie-sensitivity}에 있어요. 감쇠 계수(화살표를 따라갈 확률) 점검과 금액 기울기 점검은 엔도스랭크와 AWP의 같은 기간 지표 가족 평균을 다시 계산해요. 이 점검들은 표~\ref{tab:robustness-damping}과~\ref{tab:robustness-slope}, 그리고 그림~\ref{fig:damping-sensitivity}에 있어요. 표본 정의 단계(표~\ref{tab:robustness-sample-size})는 평가 집합에 따라 통계량이 어떻게 움직이는지 보여 줘요. 이것은 강건성(조건을 바꿔도 결과가 버티는지)의 증거가 아니에요.
\item C3: H1/RQ1 아래에 적은 순위 갈라짐이에요(표~\ref{tab:tie-sensitivity}, 그림~\ref{fig:rank-scatter}).
\item C4: 각 걷기를 자기 화살표 종류의 원시 차수와 기간 밖에서 견준 비교예요(표~\ref{tab:holdout-spenders}, \ref{tab:holdout-diff}과~\ref{tab:fresh-posthoc}). 포함 범위 점검(표~\ref{tab:holdout-receiving})과 활동 재시작을 쓰는 변형(표~\ref{tab:endorserank-restarts})도 함께 있어요.
\item C5: 고친 점수에서, 결합은 두 층을 각각 한 층만 걷는 것을 이기지 못해요. 이것은 범위가 정해진 부정적인 결과예요. 그리고 송금 걷기 하나와 견주는 규칙~B를 이 점수들에 적용하면, 그 규칙의 기준이 채워지지 않아요(표~\ref{tab:fresh-contrasts}). 등록 분석도 같은 판정에 이르렀어요(표~\ref{tab:registered-vs-revised}). 위의 RQ5와 \ref{sec:coupled-results}절을 보세요.
\end{itemize}

\dissertationSubheading[sec:map-nonclaims]{분명히 밝히는, 주장하지 않는 것들}

이 연구는 다음을 주장하지 \emph{않아요}.

\begin{itemize}
\item 어떤 점수가 일반적으로 평판을 더 잘 재는 척도라는 것, 또는 AWP가 이제 쓸모없어졌다는 것(부록~\ref{ch:appendix-er-awp}의 엔도스랭크와 AWP 비교는 구성개념을 따르고, 포함 범위에 달려 있어요).
\item 어떤 점수가 자기 화살표로 만든 대리 지표와 맞는 것이, 그 점수 바깥에서 온 증거라는 것(이것은 의도한 구성개념을 점검하는 것이고, 일부는 만든 방식 때문에 저절로 그렇게 돼요).
\item 계획의 고르게 다시 출발하기를 쓰는 엔도스랭크가 나중의 새 허락 쌍을 그냥 센 받은 허락 수가 줄 세우는 것보다 더 잘 줄 세운다는 것(두 기간 모두에서 더 못하게 줄 세웠어요).
\item 활동 재시작을 쓰는 엔도스랭크가 받은 허락 수를 이긴다는 것(그 비교는 라벨을 본 뒤에 했고, 등록한 시험을 기다리고 있어요).
\item 허락 그래프에는 여러 걸음 퍼뜨림(전파)이 있는데, 이것 때문에 엔도스랭크가 그냥 센 수를 이겼다는 것(\ref{sub:theory-closed-form}절).
\item 청산을 피하는 데에서, 그냥 센 수의 수준(R1)이나 페이지랭크 수준(R2)에서 허락 화살표가 송금 화살표보다 앞선다는 것, 또는 거기서 엔도스랭크가 AWP를 이긴다는 것(고친 점수에서도, 등록 분석에서도 어느 규칙도 채워지지 않았어요).
\item 어떤 점수가 신용도(돈을 갚을 만한 정도), 채무 불이행(빚을 못 갚는 일) 위험, 또는 거래 실력을 잰다는 것(W1에서 허락한 주인 수(허락을 해 주는 서로 다른 주인의 수)는 이익을 낸 비율과 손실이 아닌 비율에서 트레이더를 송금 수보다 못하게 줄 세웠어요. 그리고 자료에는 채무 불이행 라벨이 없어요).
\item 허락을 받는 것이 계약 트레이더(GMX V2에서 거래하는 계약)로 하여금 청산을 피하게 만든다는 것(표본이 작고, 허락을 받은 트레이더는 포지션을 더 많이 닫아요. 그리고 송금 수도 이 라벨을 줄 세워요).
\item 금액 기울기 $b=\ln 3/m$, C-PR 허락 층의 상한 $U$, 또는 엔도스랭크의 고르게 다시 출발하기가 가장 좋은 선택이라는 것(기울기와 상한은 등록 분석의 결과를 안 뒤에 수정 과정에서 정했어요. 그리고 표~\ref{tab:robustness-slope}의 금액 기울기 점검은 기울기는 바꾸지만 상한은 바꾸지 않아요. 고르게 다시 출발하기는 계획에서 정했고, 활동 재시작과의 비교는 라벨을 본 뒤에 했어요).
\item 고친 점수에서 나온 규칙~A와~B, 그리고 R1--R3의 판정이 등록한 시험이라는 것(이 규칙들은 등록 분석을 위해 정했고, 그 결과를 안 뒤에 정의한 점수에 적용했어요. 등록 분석의 판정은 부록~\ref{sec:map-registered}에 있어요).
\item C-PR이 송금만 걷는 걷기보다 낫다는 것(규칙~B는 새 송금자에서도, 트레이더에서도 실패해요), 또는 AWP보다 낫다는 것(두 기간 모두에서 새 송금자에서 AWP에 뒤져요).
\item 두 화살표 종류를 어떻게 합쳐도 한 층 방법들을 이길 수 없다는 것. 이 연구가 보인 것은 여기서 시험한 방법, 곧 $\lambda$를 고정하고 점마다 볼록 혼합(두 값을 정한 비율로 섞기)을 하는 방법이 이기지 못한다는 것뿐이에요.
\item 엔도스랭크의 더 짧은 푸는 시간이 그 알고리즘의 장점이라는 것(이것은 더 적은 화살표 수에서 나와요), 또는 코호트의 에고 네트워크(코호트 계약에 닿는 화살표만 모은 그래프)보다 큰 그래프에서 실행 시간이 어떻다는 것.
\item 엔도스랭크가 시빌 농장(공격자가 만든 가짜 주소 무리)을 견뎌 낸다는 것(평범한 점수나 중앙값보다 높은 점수에 이르는 데는 비용이 적게 들어요. 그리고 다섯 토큰 가운데 하나의 허락에는 잔고가 필요 없어요. 그래서 토큰 목록이 농장의 비용을 올리지 않아요. \ref{sec:sybil-cost-model}절).
\item 결과가, 살펴본 몇 달 동안 아비트럼 원(Arbitrum One)에서 WETH, USDC, USDT0, ARB, WBTC로 서로 다른 주인 셋 이상에게 허락을 받은 계약을 넘어서도 들어맞는다는 것. 또는 다른 토큰의 허락과 송금에도 들어맞는다는 것(수정 분석은 이것들을 빼 두었어요. 표~\ref{tab:usd-tokens}).
\item 온체인(블록체인 위의) 점수가 담보, 감사(코드 점검), 신원 시스템을 대신한다는 것. 또는 이 점수가 오프체인(블록체인 밖의) 정보를 쓴다는 것(이 연구는 그런 정보를 쓰지 않아요).
\end{itemize}

\dissertationSubheading[sec:map-artifacts]{재현 자료}

아래의 경로는 부록~\ref{ch:appendix-eval}에서 이름을 밝힌 실험 폴더를 기준으로 한 상대 경로예요. 본문의 장들이 보고하는 수정 분석은 자기 폴더가 따로 있어요. 등록 분석의 폴더는 얻었을 때 그대로 남겨 두었어요.

\begin{enumerate}
\item 모든 방법의 계약별 점수예요. 같은 기간의 코호트 점수, 그리고 W0과 W1의 스펜더 점수와 트레이더 점수가 parquet 파일에 들어 있어요. 수정 분석의 것은 \texttt{data/2-processed-tables-and-evaluations/revised-usd/}(\texttt{same-window/}, \texttt{holdout-w0/}, \texttt{window-w1/})에 있어요. 등록 분석의 것은 \texttt{data/2-processed-tables-and-evaluations/}의 \texttt{same-window/}, \texttt{holdout-w0/}, \texttt{registered-w1/}에 있어요.
\item 기계가 읽을 수 있는 요약이에요. 같은 기간(감쇠 계수를 바꿔 본 실험 포함), W0, W1, 강건성 점검, 보충 대비, 시빌 모형, 벤치마크(속도 재기), 자료 설명의 요약이에요. 여기에는 모든 계수, 구간, 대비, 판정이 들어 있어요. 수정 분석의 요약은 금액 기울기 점검과 함께 \texttt{data/3-published-results-for-thesis/revised-usd/}에 있어요. 등록 분석의 요약은 상위 토큰 점검과 함께 \texttt{data/3-published-results-for-thesis/}에 있어요. 폴더마다 자기 \texttt{results\_digest.md}와 \texttt{SHA256SUMS.txt}가 있어요.
\item 분석 계획(\texttt{docs/analysis\_plan.md})과 설정 파일이에요. 둘 다 이 연구의 로그를 하나도 뽑기 전인 2026년 10월 9일에 \texttt{43d8045}로 커밋한 그대로예요.
\item W1 등록 문서(\texttt{docs/registration\_w1.md})예요. \texttt{0d3a057}로 커밋한 그대로예요. 그리고 추출 목록 \texttt{data/2-processed-tables-and-evaluations/registered-w1/extraction\_manifest.json}도 있어요. 이 목록에는 커밋, W1 추출을 시작한 시각과 끝낸 시각, 강제 실행 옵션(override)을 쓰지 않았다는 것, 그리고 SHA-256 해시가 적혀 있어요. 등록 문서의 해시는 \texttt{2126c8dc3fe3c0f9273aa4ecfc1e15bddc1f5a6f65e98548b9528f2b18ca2a6b}이고, 설정 파일의 해시는 \texttt{d6a30adbcabcf511d6892332da73639efdbece25891b392e13e90faf09b7a833}이에요.
\item 금액 수정 문서(\texttt{docs/revision\_price\_weighting.md})예요. 2026년 10월 10일 07:54~UTC에 \texttt{a897475}로 커밋했어요. 다섯 토큰의 선택은 08:12~UTC에 \texttt{09dfc32}로, 고친 허락 상한과 상수들은 11:24~UTC에 \texttt{2705b28}로 커밋했어요. 모두 등록 분석의 결과를 안 뒤이고, 고친 점수를 하나라도 계산하기 전이에요.
\item \texttt{data/0-usd-token-list-and-prices/}에 있는 토큰 목록과 가격이에요. 여기에는 체인링크(Chainlink)의 가격 피드(블록체인에 올라오는 가격 정보) 목록을 찍어 둔 것과 검토한 후보 토큰들, 조건을 채운 토큰들, 하루 달러 가격과 그 가격이 나온 DefiLlama 데이터 점들, 상수 $m_T$, $m_A$, $U$, 그리고 이 가격을 체인에서 읽은 Chainlink 피드와 견준 결과가 있어요.
\item 비용 장부 \texttt{data/bq\_ledger.json}이에요. 여기에는 모든 빅쿼리(BigQuery) 작업과 그 매개변수, 그리고 요금이 매겨진 바이트 수가 있어요.
\item \texttt{data/1-raw-blockchain-extracts/}에 있는 뽑은 표들이에요. 계정 종류가 붙은 코호트 후보, GMX~V2 계정과 닫기 기록이 있어요. \texttt{graph-tables/}에는 등록 분석의 쌍 표, 점 표, 대리 지표, 라벨, 상위 토큰 부분 그래프가 있어요. \texttt{graph-tables-usd/}에는 다섯 토큰에 대한 수정 코호트와 그 쌍 표, 점 표, 대리 지표, 라벨이 있어요. 모두 90\,MB보다 작은 parquet 조각으로 나뉘어 있고, 조각마다 원본 표를 적은 JSON 기록이 붙어 있어요.
\item 이 논문 폴더에 있는 내보낸 표 조각(\texttt{results/tables/})과 벡터 그림(\texttt{Figures/generated/})이에요. 이것들은 수정 분석의 요약으로 만들었어요. 표~\ref{tab:registered-vs-revised}은 등록 분석의 요약도 읽어요.
\end{enumerate}

표와 그림은 뽑은 표 위에서 평가 스크립트를 돌린 다음, 같은 요약에서 표와 그림을 내보내면 다시 만들어져요(부록~\ref{sec:pipeline-stages}). \texttt{CONTRACT\_REP\_VARIANT=usd}로 두면 스크립트는 다섯 토큰의 표를 읽어요. 이 표는 \texttt{scripts/run\_usd\_pipeline.py}가 빅쿼리에서 만들어요. 그리고 스크립트는 수정한 결과를 써요. 이 변수를 설정하지 않으면 스크립트는 등록 분석의 결과를 다시 만들어요.

\dissertationSubheading[sec:map-c6]{C6: 어느 화살표 종류로도 만들지 않은 결과 값}

C6은 두 그래프 어디에도 들어 있지 않은 트레이더 결과 값을 놓고 두 화살표 종류를 견줘요.

\begin{itemize}
\item 진술: GMX~V2 계약 트레이더의 청산 없이 닫은 비율을 봤어요. 이때 W1과 함께 등록한 규칙을 고친 점수에 적용했어요. 그랬더니 허락한 주인 수가 송금 수보다 그냥 센 수의 수준에서 앞선다는 것(R1)은 W1에서 뒷받침되지 않아요. W1은 규칙이 판정을 내리는 기간이에요. W0에서는 P가 양수인데도 그래요. 페이지랭크 수준에서 앞선다는 것(R2)도 뒷받침되지 않아요. 그리고 R3이 적용돼요. 등록할 때 정한 문구로 말하면 이래요. 나중에 청산을 피하는 데에서, 그냥 센 수의 수준으로는 허락이 송금보다 앞서지 않아요. 페이지랭크 버전들은 어느 쪽이 나은지 가려지지 않았어요. 그리고 이 신호는 강제 청산을 피하는 것에 관한 것이지, 거래 성공에 관한 것이 아니에요. 이 판정들은 재분석이지 등록한 시험이 아니에요. 등록 분석도 같은 판정에 이르렀어요(표~\ref{tab:registered-vs-revised}).
\item 증거: 표~\ref{tab:neutral-label-contrasts}과~\ref{tab:neutral-label-recipients}(\ref{sub:neutral-label-results}절)에 있어요. 두 기간의 트레이더 라벨 여섯 개 모두는 표~\ref{tab:neutral-label-grid}에 있어요. 절차는 \ref{sub:neutral-label-protocol}절에 있어요. 등록 내용은 \texttt{docs/registration\_w1.md}의 3절에 있어요. 이 문서는 스캔~B 전에 \texttt{0d3a057}로 커밋했어요. 이 등록 내용을 고친 점수에 적용한 것은 \texttt{docs/revision\_price\_weighting.md}에 있어요(\ref{sec:amount-revision}절).
\item 주의할 점: 트레이더 코호트는 작아요. W1에서는 계약 $70$개, W0에서는 $88$개예요. 2026년에 계약 계정이 GMX~V2에서 거래를 적게 했기 때문이에요(표~\ref{tab:gmx-quarters}). 다섯 토큰 가운데 하나의 허락을 받은 적이 있는 트레이더는 두 기간 모두에서 전부 청산을 피했어요. 다른 트레이더는 약 네 개 중 세 개꼴로 피했어요. 하지만 허락을 받은 트레이더는 포지션을 더 많이 닫았어요. 그리고 송금 수도 이 라벨을 줄 세워요. 모든 트레이더가 계약이에요. 그래서 계정 종류는 고정되어 있어요. 이 점수를 걸러 내는 데 쓰려면, 먼저 더 큰 표본에서 등록한 시험을 해야 해요(\ref{sec:future-work}절).
\end{itemize}

\dissertationSubheading[sec:map-registered]{등록 분석}

분석 계획과 W1 등록 문서는 금액을 토큰의 원래 단위(소수점 조정을 하지 않은 가장 작은 단위)로 재고, 모든 ERC-20(토큰 공통 규칙) 토큰을 쓰고, 값 변환(금액이 아무리 커도 1을 넘지 않게 바꾸는 식)에 $b=1$을 쓰기로 정했어요. 금액은 그 분석의 결과를 안 뒤에 고쳤어요(\ref{sec:amount-revision}절). 본문의 장들은 수정 분석을 보고해요. 등록 분석은 얻은 그대로 여기에 남겨 두었어요. 표~\ref{tab:registered-vs-revised}은 등록 분석의 판정 대비와 판정을, 고친 점수에서 구한 같은 대비 옆에 나란히 놓아요. 이 표의 마지막 네 행은 두 기간의 차수 대비예요. 이 대비들은 두 분석 모두에서 라벨을 본 뒤에 계산했고, 판정이 없어요. 등록 열에 있는 모든 판정은 이 연구가 그 자료를 뽑기 전에 정한 시험의 결과예요. 수정 열에 있는 판정은 어느 것도 그렇지 않아요. 규칙을 등록 분석의 결과를 안 뒤에 정의한 점수에 적용했기 때문이에요. 모든 계수, 구간, 판정을 담은 등록 분석의 전체 결과는 계산한 그대로 \texttt{data/3-published-results-for-thesis/}(\texttt{results\_digest.md})에 공개했어요. 수정한 결과는 그 아래 폴더 \texttt{revised-usd/}에 있어요.

\IfFileExists{results/tables/registered-vs-revised.tex}{% Generated by experiments/scripts/export_latex.py -- do not edit by hand.
% Source: experiments/data/2-processed-tables-and-evaluations/same-window/eval_summary.json
% Source: experiments/data/2-processed-tables-and-evaluations/holdout-w0/eval_summary.json
% Source: experiments/data/2-processed-tables-and-evaluations/registered-w1/eval_summary.json
% Source: experiments/data/2-processed-tables-and-evaluations/revised-usd/same-window/eval_summary.json
% Source: experiments/data/2-processed-tables-and-evaluations/revised-usd/holdout-w0/eval_summary.json
% Source: experiments/data/2-processed-tables-and-evaluations/revised-usd/window-w1/eval_summary.json
\begin{table}[htbp]
\centering
\caption[등록 분석과 수정 분석의 결정 대비.]{등록 분석의 결정 대비와 판정이에요(금액은 원래 단위, 모든 ERC-20 토큰, $b=1$). 분석 계획과 W1 등록 문서에 따라 얻은 그대로예요. 그 옆에 수정 분석의 점수에서 구한 같은 대비를 놓았어요(금액은 달러(USD), 목록에 오른 토큰만). 수정은 등록 분석의 결과를 본 뒤에 정했어요. 수정 분석 열은 등록한 규칙과 기준값을 수정 분석의 점수에 적용한 것이고, 등록한 검사가 아니에요. 마지막 묶음에는 두 분석의 어떤 규칙도 판정에 쓰지 않는 차수 대비가 있어요. 이 대비들은 두 분석 모두에서 보고해요. $\Delta\tau$와 95\% 짝지은 부트스트랩 구간을 적었어요(다시 뽑기 400번과 2{,}000번, 시드 42). 규칙과 그 기준값은 \ref{sec:holdout-protocol}절과 \ref{sub:fresh-holdout}절에 있어요.}
\label{tab:registered-vs-revised}
\fitwidth{%
\begin{tabular}{lcc}
\toprule
항목 & 등록 분석 & 수정 분석 \\
\midrule
매칭 코호트 (계약 수) & 27{,}844 & 15{,}107 \\
W0 스펜더 코호트 & 33{,}101 & 17{,}050 \\
W1 스펜더 코호트 & 34{,}587 & 18{,}152 \\
W1 트레이더 & 70 & 70 \\
EndorseRank 대 AWP, 같은 기간 ($\tau$) & 0.217 [0.209, 0.225] & 0.174 [0.163, 0.185] \\
\midrule
A, W0: C-PR 빼기 C-PR ($\lambda=1$), 새 허락 쌍 & $+0.008$ [0.003, 0.013] & $+0.029$ [0.022, 0.036] \\
A, W0: C-PR 빼기 C-PR ($\lambda=0$), 새 송금자 & $-0.030$ [-0.034, -0.026] & $-0.034$ [-0.040, -0.029] \\
규칙 A, 주된 부분 & 충족 안 됨 & 충족 안 됨 \\
규칙 A, 보조 부분 & 충족 안 됨 & 충족 안 됨 \\
\midrule
B, F1: C-PR 빼기 C-PR ($\lambda=0$), 새 허락 쌍 & $+0.101$ [0.092, 0.111] & $+0.110$ [0.098, 0.122] \\
B, F2: C-PR 빼기 C-PR ($\lambda=0$), 새 송금자 & $-0.025$ [-0.029, -0.021] & $-0.028$ [-0.033, -0.023] \\
B, F3: C-PR 빼기 C-PR ($\lambda=0$), 청산 없이 닫은 비율 & $-0.044$ [-0.100, -0.002] & $+0.000$ [-0.014, 0.010] \\
규칙 B (F1, F2, F3) & 충족 안 됨 & 충족 안 됨 \\
규칙 B, 민감도 1 (AWP를 비교 대상으로 함) & 충족 안 됨 & 충족 안 됨 \\
규칙 B, 민감도 2 (외톨이 점) & 충족 안 됨 & 충족 안 됨 \\
F6 (규칙 A, 주된 부분, W1에서) & 불성립 & 불성립 \\
\midrule
P, W1: 받은 허락 수 빼기 송금 입차수 & $+0.063$ [-0.026, 0.148] & $+0.084$ [-0.013, 0.176] \\
Q, W1: EndorseRank (활동에 따른 재시작) 빼기 AWP & $+0.051$ [-0.085, 0.213] & $-0.041$ [-0.135, 0.050] \\
R1 (개수 수준의 앞섬) & 충족 안 됨 & 충족 안 됨 \\
R2 (페이지랭크 수준의 앞섬) & 충족 안 됨 & 충족 안 됨 \\
R3 (특이성) & 적용됨 & 적용됨 \\
\midrule
W0: EndorseRank 빼기 받은 허락 수, 새 쌍 & $-0.011$ [-0.016, -0.006] & $-0.019$ [-0.026, -0.011] \\
W0: AWP 빼기 송금 입차수, 새 송금자 & $+0.038$ [0.035, 0.042] & $+0.043$ [0.038, 0.048] \\
W1: EndorseRank 빼기 받은 허락 수, 새 쌍 & $-0.029$ [-0.032, -0.025] & $-0.034$ [-0.038, -0.029] \\
W1: AWP 빼기 송금 입차수, 새 송금자 & $+0.020$ [0.016, 0.023] & $+0.022$ [0.017, 0.027] \\
\bottomrule
\end{tabular}
}
\end{table}
}{}
